\scanpage{40}
% archive sheet number: 39

De plus, (1) (2) (3) [\unsure{impliquent|implique|l'idée que} l'existence d'un \(G\) \deleted{\unsure{de}}\unsure{simple} sur \(S\) donnant \(\mathfrak g\)] prouvent aussi [pour \(G\) connexe] que \(G_{\unsure{\bullet|0}}\) n'a pas de variations infinitésimales non triviales.

Question de relèvement de la structure \(G\)\inserted{ (\(G\) connexe)}, équivaut à la question de relèvement de la structure \(\mathfrak g\). On trouve une obstruction dans \(H^3(\mathfrak g,\mathfrak g)\). Donc

\deleted{\(H^{3}\)}

(3) Si \(H^3(\mathfrak g,\mathfrak g)=0\), alors on peut relever de façon \unsure{essentiellt|essentiellement} unique \(\mathfrak g\), ou \(G\) de façon équivalente.

\uline{Question} Définir \(H^{\unsure{*|\alpha}}(G,\mathfrak g)\to H^*(\mathfrak g,\mathfrak g)\)

[plus généralement, si on a une représentation linéaire de \(G\) \ldots]. Suite spectrale reliant \(H^*(G,\ )\) et \(H^*(\mathfrak g,-)\) ?? Elle exprimerait la \textquotedblleft différence\textquotedblright\ entre le global, et l'infinitésimal.

Supposons \(G\) semi-simple adjoint. A-t-on alors

\[
H^*(\unsure{G|h},\mathfrak g)=H^*(\mathfrak g,\mathfrak g)=0\qquad ??
\]
